% Generated draft: requires mathematical and source review.
\documentclass[11pt]{article}
\usepackage[margin=25mm]{geometry}
\usepackage{fontspec}
\setmainfont{Latin Modern Roman}
\usepackage{amsmath,amssymb,mathtools,unicode-math}
\setmathfont{Latin Modern Math}
\usepackage{enumitem,needspace,xurl,hyperref}
\setlength{\parindent}{0pt}
\setlength{\parskip}{5pt}
\setlength{\emergencystretch}{3em}
\newcommand{\N}{\mathbb{N}}
\newcommand{\Z}{\mathbb{Z}}
\newcommand{\Q}{\mathbb{Q}}
\newcommand{\R}{\mathbb{R}}
\newcommand{\C}{\mathbb{C}}
\newcommand{\F}{\mathbb{F}}
\newcommand{\A}{\mathbb{A}}
\newcommand{\PP}{\mathbb{P}}
\newcommand{\E}{\mathbb{E}}
\newcommand{\eps}{\varepsilon}
\newcommand{\dd}{\,\mathrm{d}}
\begin{document}
\section*{UnsolvedMath ID: 30006038}
\textbf{Source problem code:} catalogue entry
\subsection{Definitions and mathematical statement}
\textbf{Unresolved mathematical statement.} The supplied record is a workshop-summary fragment covering several talks, not an identifiable standalone mathematical problem. Its supplied statement status is \texttt{unrecoverable}, and no clean statement is supplied. Consequently, no precise mathematical assertion or associated system of definitions can be recovered without adding unsupported hypotheses, domains, boundary conditions, decay rates, or quantifiers. The following preserves the mathematical context and its attribution; it does not promote the reported results or suggested directions to an open problem.

\begin{itemize}
\item Sung-Jin Oh described a newly developed method for generating initial data sets with good localisation and fall-off properties, described in the passage as flexible and promising and as generalising and complementing existing gluing results. No constraint equations, class of initial data, localisation conditions, or fall-off rates are specified.
\item Anna Sakovich described recent results concerning the definition and geometry of the mass aspect function for asymptotically hyperbolic, rather than asymptotically Euclidean, initial data sets. The passage suggests studying such data under very weak decay assumptions, with possible relevance to stability of geometric inequalities such as the positive mass theorem. It gives neither precise decay assumptions nor a stability assertion.
\item Markus Wolff described a geometric flow along null hypersurfaces that allows construction of foliations by surfaces of constant spacetime mean curvature in null cones near null infinity. The passage relates this notion to foliations considered by Carla Cederbaum and Anna Sakovich in initial data sets. It supplies no definition of the curvature, flow equation, spacetime hypotheses, or meaning of the neighbourhood of null infinity.
\item Annegret Burtscher described a characterisation of global hyperbolicity via completeness of the null distance, presented as a Lorentzian analogue of the Hopf--Rinow theorem. The passage does not specify the construction of the null distance or the hypotheses needed for the characterisation.
\item James Lucietti described results fully characterising the near-horizon geometry of stationary spacetimes hosting a degenerate Killing horizon, together with their use in proving rigidity. The spacetime dimension, field equations, regularity, and other hypotheses are not given.
\item Jan Sbierski described a result showing Lipschitz inextendibility of weak null singularities, which the passage says are expected to form in the interior of generic rotating black holes. It states that such an inextendibility statement would be necessary to deduce versions of Penrose's strong cosmic censorship conjecture. Neither the relevant class of spacetimes, the meaning of genericity, the extension category, nor a particular version of the conjecture is specified.
\end{itemize}

The supplied label \emph{open} is a registry normalisation of an unresolved extraction, not a verified assertion that any one of these reported results or directions is an open problem. No target among them is selected by the source record.
\subsection{Short English statement}
This record summarises several talks on mathematical general relativity and supplies no standalone problem. A precise statement cannot be reconstructed without choosing a target and adding missing hypotheses. The supplied \emph{open} label reflects unresolved source extraction; this statement-only draft requires human mathematical and source review.
\subsection{Sources}
\begin{itemize}
\item Supplied local UnsolvedMath registry record: \texttt{data/unsolvedmath/problems.json}, permanent ID \(30006038\), title \emph{Workshop Summary on Mathematical General Relativity}, exact source problem code \texttt{catalogue entry}. The supplied set is \emph{Oberwolfach Reports Open Problems}; the proposed year is 2024, and no proposer is identified.
\item Supplied source citation: \emph{Mathematical Aspects of General Relativity (2025)}, Oberwolfach Reports, pp.\ 2033--2112, DOI \texttt{10.4171/owr/2024/36}. Supplied source and background link: \url{https://doi.org/10.4171/owr/2024/36}. The supplied literature-source label is \emph{Oberwolfach Report 36 (2024)}, with the same link: \url{https://doi.org/10.4171/owr/2024/36}. The differing year labels are retained without reconciliation.
\item Supplied literature assessment, checked on \texttt{2026-08-22}: the extraction is contextual report prose rather than an identifiable mathematical assertion; its unresolved status was normalised to \emph{open}. This is supplied dated curation, not an independently verified literature assessment. The background expressly requests independent verification of time-sensitive status and cited claims.
\item Supplied statement-verification note: the official report is said to confirm that the extraction is workshop-summary prose covering several talks and contains no standalone problem. The note says that no separate clean statement was promoted and that the original record remains publicly visible and retained for provenance. These are supplied verification claims; no source visit or independent verification is asserted here.
\item Mathematical attribution in the supplied passage is to Sung-Jin Oh, Anna Sakovich, Markus Wolff, Annegret Burtscher, James Lucietti, and Jan Sbierski, with additional references to Carla Cederbaum and Anna Sakovich for the foliations and to Penrose for strong cosmic censorship. These attributions do not identify a proposer of a standalone problem.
\item The passage records acknowledgements by the MFO and workshop organisers to the National Science Foundation for supporting junior researchers through grant \texttt{DMS-2230648}, ``US Junior Oberwolfach Fellows'', and to the Simons Foundation for supporting Todd Oliynyk and Igor Rodnianski through the ``Simons Visiting Professors'' programme at the MFO.
\end{itemize}
\end{document}
